\documentclass[11pt]{article}

\usepackage[margin=1in]{geometry}
\usepackage{amsmath}
\usepackage{amssymb}
\usepackage{booktabs}
\usepackage{array}
\usepackage{longtable}
\usepackage[hidelinks]{hyperref}
\usepackage{microtype}
\usepackage{xcolor}
\usepackage{listings}

\lstset{basicstyle=\ttfamily\small, columns=fullflexible, keepspaces=true, breaklines=true}

\newcommand{\code}[1]{\texttt{#1}}

\title{Formal Equivalence of Hazard3 and the RISC-V Sail Model\\
\large An unbounded, open-tool proof of an RV32IMC configuration of the RP2350 core,
modulo one stated axiom, and the six divergences it found}
\author{Leo Strijbos}
\date{September 2026}

\begin{document}
\maketitle

\begin{abstract}
Hazard3 is the RISC-V core in the Raspberry Pi RP2350. We port the lowRISC/Oxford
Ibex proof harness to Hazard3 and prove, with Yosys and the rIC3 model checker,
that an RV32IMC + Zicsr machine-mode-only build of Hazard3 v1.1.1 is trace
equivalent to the RISC-V Sail specification: at every specification query the
program counter, the six checked CSRs and the register operands agree, and the
core's memory requests are the specification's memory requests, for every
instruction sequence and without a depth bound. All 18 top-level goals close.
The proof rests on 316 helper lemmas in 108 layers and on one axiom, the
arithmetic result of the sequential multiplier/divider, which the Ibex proof
also assumes. Six divergences had to be fixed before the proof could close:
five in the Hazard3 RTL (reserved \code{C.ADDI4SPN} encodings execute,
\code{mstatus.MPRV} writable without U-mode, a trapping \code{meicontext} write
with a side effect, \code{FENCE} trapping on non-zero rs1/rd, and vectored trap
entry using OR rather than addition) and one in the lowRISC revision of the Sail
model (\code{mstatus.TW} not masked without U-mode). All five RTL divergences
reproduce on Hazard3's own simulator; one had been fixed independently on the
development branch before we found it. The lemma discovery and the proof runs
were done by an LLM agent under human direction; we report what that required.
\end{abstract}

\section{Introduction}

A processor is conformant when it does what the ISA says for every instruction
sequence. Testing samples that space; a formal equivalence proof against an
executable specification covers it. The RISC-V Sail model is such a
specification, and the lowRISC and Oxford proof of the Ibex core against it
\cite{ibexformal} showed that an unbounded equivalence proof of a small RISC-V
pipeline is within reach of a small team with open tools.

This report applies that method to Hazard3 \cite{hazard3}, Luke Wren's
three-stage RV32 core, which is the RISC-V option in the Raspberry Pi RP2350.
Hazard3 differs from Ibex in ways that matter to a proof: an AHB5 bus with a
single data port, a compressed-instruction buffer with a lock that pins the
decode stage across a predicted branch, a sequential 32-cycle multiplier and
divider, and an interrupt controller of its own. None of the Ibex lemmas
transfer as written. What transfers is the shape of the proof and the
discipline around it.

\subsection{What is claimed}

Let $S$ be the specification state, $A$ the micro-architectural state of
Hazard3, $\mathit{abs} : A \to S$ the abstraction that reads the architectural
registers out of the pipeline, and let $\mathit{Spec}$ and $H3$ be the one-step
functions of the specification and of the core, with $\mathit{SpecOut}$ and
$H3\mathit{Out}$ their memory operations. The 18 goals establish two things.

\begin{itemize}
\item \textbf{Wrap}: $\mathit{Spec}(\mathit{abs}(a), i) = \mathit{abs}(H3(a, i))$
  for every reachable $a$ and instruction $i$. One step of the specification and
  one step of Hazard3 land in the same architectural state.
\item \textbf{Load, Store, NoMem, Mem}: $\mathit{SpecOut}(\mathit{abs}(a), i) =
  H3\mathit{Out}(a, i)$. The memory operation Hazard3 performs is the one the
  specification performs, with the same address, direction, byte lanes and data.
\end{itemize}

By induction over queries, the $n$th memory operation agrees for every $n$, and
the architectural state agrees at every query, from a common start. The claim is
conditional on instructions continuing to retire: liveness is not proven, the
same caveat the default Ibex flow carries.

The proof is of a patched copy of the RTL. Five one-line changes fix the five
Hazard3 divergences of section~\ref{sec:findings}; without them, the
\code{Wrap} goals for mstatus, mcause, mepc and the PC are false. The specification is the lowRISC revision of the
Sail model with three patches of ours, one of which fixes a defect in that
revision. Every claim in this report is therefore about Hazard3 v1.1.1 with
those five changes, in the configuration of section~\ref{sec:dut}, under the
assumptions of section~\ref{sec:scope}.

\subsection{Summary of results}

\begin{itemize}
\item All 18 goals close, each with the verdict \emph{proven modulo axioms
  [CtiWbMulDiv]} from a chain checker that walks the proof log and requires
  every constraint used at proof time to be itself proven on the same model.
\item One axiom, stated in full in the proof file: for the eight RV32M
  instructions, the value the datapath produces equals the specification's
  value. Everything around it (that the instruction retires, writes only rd,
  does so at the right time, does not trap, advances the PC) is proven.
\item Six divergences, each a machine-checked counterexample found by bounded
  model checking from reset, each reproduced independently of the proof on
  Hazard3's CXXRTL testbench with a program of a few dozen hand-encoded
  instructions.
\item The lemma loop ran from 17 to 20 September 2026 on one 48-core machine,
  in 44 rounds, and tried 258 distinct lemma candidates, of which 14 were
  refuted for real.
\end{itemize}

\section{Setup}

\subsection{The device under test}
\label{sec:dut}

Hazard3 at tag v1.1.1, commit \code{8af99293}, the head of its \code{stable}
branch. The configuration is Hazard3's minimum that still has a multiplier and
compressed instructions: RV32IMC + Zicsr, \code{MULDIV\_UNROLL = 1}
(the sequential shift-add multiplier and restoring divider, not the fast
multiplier), no A, no bit-manipulation, no Zcb/Zcmp, no Zifencei, no Xh3
extensions, machine mode only, no PMP, no debug, no branch predictor, one
external interrupt line plus the software and timer interrupts. The RP2350
ships a larger configuration (RV32IMAC with Zba, Zbb, Zbs, Zbkb, Zcb, Zcmp, the
Xh3 extensions, U-mode, PMP and debug). What is proven here is a subset of
that core's behaviour, not the RP2350's build.

The RTL is read by Yosys through yosys-slang. Two kinds of patch are applied to
a copy, and a manifest is written with every build. \emph{Tool patches} change
no behaviour (two functions are marked \code{automatic} so the front end accepts
them). \emph{Conformance patches}, applied only with \code{H3\_CONFORMANCE=1},
change behaviour and are the fixes for the five findings. A proof of the
tool-patched design is a proof of the shipped design; a proof of the
conformance-patched design is a proof that the fixes close the divergences.

\subsection{The specification}

The Sail model is lowRISC's fork of \code{riscv/sail-riscv} at commit
\code{81a266b6}, the revision Ibex's formal flow pins, compiled to
SystemVerilog by lowRISC's Sail back end. Three patches of ours are applied.

\begin{itemize}
\item \textbf{Gate 6} undoes two Ibex specialisations that are not
  architectural: lowRISC had replaced upstream's WARL legalisation of
  \code{mtvec} with a 256-byte-aligned, vectored-only policy and added an
  \code{mcause} legalisation carrying Ibex's internal interrupt encoding; and
  its memory rework had deleted the call to \code{check\_misaligned}, so the
  model split every misaligned access regardless of the platform flag. Both
  are restored to upstream's behaviour, parameterised so that the Ibex policy
  is one setting of the parameter.
\item \textbf{Gate 7} makes \code{FENCE.I} decode conditional on Zifencei,
  which this configuration of Hazard3 does not implement.
\item \textbf{Gate 8} masks \code{mstatus.TW} when U-mode is absent. This is
  finding 6: a defect in the lowRISC revision, fixed upstream in February 2025.
\end{itemize}

\subsection{Tools}

Yosys with yosys-slang for the RTL, lowRISC's Sail back end for the
specification, \code{psgen} for the proof script, \code{aig-manip} (lowRISC,
Rust) to select a property and its constraint set out of the AIGER model, and
rIC3 for bounded model checking, k-induction and IC3. All come from the Nix
flake of the Ibex formal flow at Ibex commit \code{405c6d1d}. The proof
machine was one AWS r7i.12xlarge (48 vCPUs, 384 GB), running 18 proofs in
parallel.

The full model, \code{all.aig}, has 10.4 million AND gates and 4762 latches;
the specification appears as a combinational block with 1.8 million free
inputs. A second model, \code{core.aig}, is the same harness with the
specification's outputs left as free inputs. It is built from the identical
property set, and is the basis of the pre-screen of section~\ref{sec:twomodel}.

\subsection{The harness}

The harness is a port of Ibex's \code{dv/formal}. Its parts, in the order they
were built and validated (the bring-up notes in \code{docs/BRINGUP.md} record
each gate with its measurements):

\begin{enumerate}
\item \textbf{Bus environment.} Hazard3's own AHB5 slave assumptions with
  unbounded stalls, and its AHB5 master properties as proof obligations. A
  negative control (a slave that violates the protocol) was checked to be
  detected.
\item \textbf{Free state and abstraction.} Registers that reset to unknown
  values become free inputs; the abstraction function reads the six checked
  CSRs, the PC and the register file out of the pipeline at the commit point.
\item \textbf{The commit point.} The cycle at which an instruction's
  architectural effect is fixed, with the proof that every architectural change
  (PC, CSRs, register file) happens at a commit and has a cause.
\item \textbf{The follower and the ghost.} The follower is a small record,
  carried alongside the pipeline, of what the specification was told at the
  last query and what it answered. The ghost is the specification's predicted
  next state, kept until the next query and dropped when the harness declines
  to compare (bus faults, CSRs outside the allowlist). Ibex's obligations on
  the specification's answer become implications guarded on the follower.
\item \textbf{The goals.} The 18 goals of table~\ref{tab:goals}, stated so
  that each fails, if it fails, for a reason a counterexample shows. Negative
  controls prove they are not vacuous: a query does occur with the ghost valid,
  an instruction does read a register the ghost knows, memory requests do
  happen and do not always happen.
\item \textbf{The lemma loop.} Section~\ref{sec:method}.
\end{enumerate}

Every property in the proof script mentions only signals inside the core and on
the bus. The specification enters only through the follower and ghost, which is
what allows the same property set to be built into both models.

\section{Method}
\label{sec:method}

\subsection{The lemma staircase}

Each property in the proof script sits in a layer. A property is proven by
k-induction with every property in the layers beneath it as a constraint, and
nothing above or beside it. \code{aig-manip select} extracts, from the full
model, the property and exactly that constraint set; rIC3 proves it or returns
a counterexample to induction. The final proof has 108 layers, 316 helper
lemmas, 18 goals and 1 axiom. The chain for \code{WrapPc}, the last goal to
close, uses 266 constraints, 91 of them structural lemmas from the bring-up
and the rest discovered in the loop.

Lemmas in the same layer cannot see each other. This bit twice: a pair of
lemmas that each needed the other would time out together and prove instantly
once split across two layers. The runner now inserts every candidate in its own
layer unless told otherwise.

\subsection{Counterexamples to induction as the source of lemmas}

k-induction fails on a state that satisfies all constraints and the property at
step $k$ and violates it at step $k+1$. That state is usually unreachable, and
the lemma to write is the fact that rules it out. The runner produces for every
candidate, in parallel with the proof attempt, a counterexample to induction at
$k=1$ (\code{aigicti.py} builds the inductive-step model, rIC3 finds the
trace, \code{aig-manip simulate} turns it into a VCD), and a report listing
which harness predicates are false in the final frame and which pipeline
signals changed between frames. Reading that report is the whole of lemma
discovery. A typical cycle: the report shows the decode stage holding an
instruction with a register operand read from a stale port while the pipeline
is locked; the lemma says that under a lock the operand comes from the
forwarding network; that lemma proves in 200 s; the target proves in the next
round.

The $k=1$ counterexample is not a predictor of failure. Sixteen candidates with
a $k=1$ counterexample were proven by k-induction at depth 1 in the same round,
because rIC3's k-induction strengthens with the constraints in ways the bare
inductive step does not. Early versions of the runner skipped the proof when a
counterexample existed and lost those sixteen; the current one always runs
both.

\subsection{Two models}
\label{sec:twomodel}

Roughly half of all lemmas are about the pipeline alone and never mention the
follower. Those prove on \code{core.aig}, where the specification is free, and
a proof there transfers to \code{all.aig} because the models differ only in
what drives the specification's outputs. Every candidate is first tried on the
core model (bounded check to depth 15, then k-induction with a 400 s cap); a
proof is recorded and the candidate skips the full model, a refutation on the
core model is disregarded (the free specification can refute anything that
depends on it), and a timeout decides nothing. On the full model each candidate
gets a bounded check and then k-induction with a 600 s cap. The caps came from
the data: the median proof took 164 s, the 90th percentile 265 s, the longest
462 s, and a candidate that had not proven by then ran to whatever cap it was
given. The full-model cap was raised from 400 to 600 s after the 462 s proof.

Of the structural lemmas from the bring-up, 48 do not prove by k-induction on
the full model at any cap tried; after each model change they were re-proven
with IC3 (1800 s budget, all within it) instead.

\subsection{Chain of custody}

The proof script and the log together define what is proven. A lemma is
proven at a moment, against the constraint set present at that moment; if a
constraint beneath it is later refuted, removed or found never to have been
proven, the lemma is stale. On 18 September two goals were announced proven
and retracted the same day. One rested on a lemma that had only proven because
of an unproven lemma beneath it. The other had sixteen leftover
instrumentation properties as constraints, five of them vacuity controls
written to be false (``an ADD never retires''), so every proof above them had
assumed five instruction classes never execute. Bounded model checking cannot
flag a false constraint; it silently shrinks the explored space.

The chain checker, \code{tools/chain\_check3.py}, was written that day. For a
goal it reconstructs, from the insert and remove events in the log, the
constraint set present at the goal's proof line; checks that each constraint
is still present and has itself a valid proof, recursively; treats the
structural lemmas as taken on the proof header's word and says so; and prints
\emph{PROVEN}, \emph{PROVEN MODULO AXIOMS [names]} if an assumed lemma is in
the chain, or \emph{STALE} with the offending constraint. Two weaker checkers
(present and has an UNSAT line; a count of constraints) both passed states
that the third called stale. No goal has been reported since without that
output, and the rule the work ran under was that no report of a goal is made
without it.

Two consequences of the discipline. Removing a lemma means re-proving every
survivor above the lowest removed layer, which is what made stripping a bad
lemma expensive and made the runner's strip step re-prove only survivors, not
everything. And a model change (a conformance patch) means rotating the log
and re-proving every lemma from nothing; it was done twice, on 19 September
(four RTL fixes and the Sail fix) and 20 September (the \code{mtvec} fix),
each time as a full sweep with nothing stripped, followed by the chain check
on all 18 goals.

\subsection{The runner}

\code{tools/h3round\_v41.sh} runs one round: for each candidate lemma in a
list, insert it in the proof script at a stated layer, rebuild the select file,
try the core model, then the full model with the counterexample analysis in
parallel, and append \code{RESULT} lines to the log. Candidates can be tagged
\code{.dut} (skip the core model; the lemma mentions the follower) or
\code{.axiom} (insert, never attempt, log \code{AXIOM name assumed}). Goals are
tried after the candidates: a bounded check to depth 15, and k-induction only
if no counterexample to induction exists. A round of 5 to 15 candidates took
15 to 70 minutes. Nothing in the runner kills a running proof.

\section{Results}

\subsection{Goals}

\begin{table}[h]
\centering
\small
\begin{tabular}{p{0.42\textwidth}p{0.5\textwidth}}
\toprule
goal & statement \\
\midrule
WrapPc & the PC agrees at each query \\
WrapMstatus, WrapMie, WrapMcause, WrapMtvec, WrapMscratch, WrapMepc & the six checked CSRs agree at each query \\
WrapRegA, WrapRegB & both register read operands agree at each query \\
LoadEn, LoadAddr & a load is issued when and where the spec loads \\
StoreEn, StoreAddr, StoreBe, StoreWData & a store is issued when, where, which lanes, which data \\
NoMem, MemWe & no request without a spec memory op; direction agrees \\
MemNoSecondGranule & no access is split into a second granule \\
\bottomrule
\end{tabular}
\caption{The 18 goals. Every verdict on the final model is \emph{proven modulo
axioms [CtiWbMulDiv]}. \code{MemNoSecondGranule} closed first, on 18
September; \code{WrapPc} closed last, on 20 September, and the sweep on the
fully patched model re-proved everything and re-checked all 18 chains the same
day (\code{results/chain/chain\_all.txt}).}
\label{tab:goals}
\end{table}

\subsection{The axiom}

\code{CtiWbMulDiv}: whenever the M stage holds an instruction whose follower
record is clean (valid, not excepting, not a fetch fault, one the specification
could execute) and whose word is an RV32M operation (opcode OP, funct7
\code{0000001}: MUL, MULH, MULHSU, MULHU, DIV, DIVU, REM, REMU) and the
specification writes rd, then the value the core holds for that write over its
M residency equals the specification's value.

What is assumed is only the arithmetic of the iterative unit for those eight
instructions. What is proven above it is that such an instruction retires,
writes rd and only rd, at the right time, does not trap, advances the PC, and
that every other instruction class is value-correct: 39 per-class lemmas
(\code{CtiWbAdd} through \code{CtiWbCXor}), the load lemma, the CSR lemma and
their union \code{CtiWbAll}.

We tried to discharge it. A bounded check of the multiplier/divider unit
alone against a plain 64-bit product and 32-bit quotient found no
counterexample to depth 34 for MUL, MULHU and DIVU and stalled there (30
minutes per depth); k-induction and IC3 do not close a 32-cycle iterative
multiplier at the bit level. The algebraic route (AMulet-style circuit
verification against a polynomial specification) was estimated at two to three
days for the multiply half only, with division still assumed, and was not
taken. The Ibex proof carries the same assumption as \code{MType}. Hazard3's
own riscv-formal checks cover the unit in bounded form upstream.

\subsection{Statistics}

\begin{table}[h]
\centering
\small
\begin{tabular}{lr}
\toprule
helper lemmas in the final proof & 316 \\
\quad of which found in the loop (named \code{Cti*}) & 175 \\
\quad of which structural, from the bring-up & 141 \\
layers & 108 \\
constraints in the \code{WrapPc} chain & 266 \\
rounds (17 to 20 September) & 44 \\
distinct lemma candidates attempted & 258 \\
refuted for real (the candidate was false) & 14 \\
candidates that timed out at least once & 110 \\
k-induction proof time on the full model, median / p90 / max & 164 / 265 / 462 s \\
model changes with full re-proof & 2 \\
\bottomrule
\end{tabular}
\caption{Proof-run statistics from the three log files.}
\end{table}

\section{Findings}
\label{sec:findings}

Each finding started as a bounded-model-checking refutation of a goal or lemma
from reset, was traced to an RTL or Sail line, was fixed by a conformance
patch, and later was reproduced outside the proof with a program of a few
dozen instructions on Hazard3's own CXXRTL testbench at v1.1.1
(\code{findings/upstream/repro/}). Each program exits 1 on v1.1.1 and 0 on
v1.1.1 with the patch applied; the same programs were run on the \code{develop}
head of 29 August 2026 (\code{ba0c83c6}). Hazard3's contribution guide asks
that reports contain no RTL suggestions, so the drafts describe behaviour only.

\paragraph{1. Reserved C.ADDI4SPN encodings execute.}
\code{C.ADDI4SPN} has $\mathit{nzuimm} = \mathit{instr}[12{:}5]$ and the
encodings with $\mathit{nzuimm} = 0$ are reserved. The decompressor tests
\code{instr[12:2]}, which also requires $rd' = 0$, so only the all-zero
halfword traps; the other seven encodings (\code{0x0004} to \code{0x001c})
retire as \code{addi rd', sp, 0}. Every sibling reserved test in the file
checks exactly the field the specification names. Present on \code{develop}.

\paragraph{2. mstatus.MPRV writable without U-mode.}
With \code{U\_MODE = 0} the privileged specification makes MPRV read-only 0.
The mstatus write arm gates MPP and TW on \code{U\_MODE} but not MPRV, which
sets and reads back (mstatus \code{0x00021800} after \code{csrs mstatus,
0x20000}). The bit feeds the effective load/store privilege. Present on
\code{develop}.

\paragraph{3. A trapping meicontext write modifies mie.}
With \code{EXTENSION\_XH3IRQ = 0} the CSR does not decode and a write raises
illegal-instruction, but the arm that applies the write's side effect to
\code{mie} is gated only on the address, so \code{mie} goes from 0 to 8 in the
trap-entry cycle. An instruction that raises an exception must not change
architectural state. Present on \code{develop}.

\paragraph{4. FENCE with non-zero rs1 or rd traps.}
The opcode pattern pinned rs1 and rd to zero; the unprivileged specification
says base implementations shall ignore those fields. Widening the pattern
alone would make \code{fence} with rd = x4 write x4, so the decode arm must
also force rd to x0. The maintainer fixed this on \code{develop} on 22 August
2026 (\code{6adda670}) from a newer riscv-arch-test run, before this work found
it on 19 September; the fix is not on \code{stable}. We record it as an
independent rediscovery and do not file it.

\paragraph{5. Vectored trap entry uses OR.}
The specification sets $pc = \mathit{BASE} + 4 \cdot \mathit{cause}$; Hazard3
computes $\mathit{BASE} \mid 4 \cdot \mathit{cause}$ and documents that the
table must be aligned to its size. That is a legitimate implementation choice
if the register enforces it, but with the default \code{MTVEC\_WMASK} every
BASE bit is writable, so an unaligned BASE in vectored mode is stored, reads
back, and sends a software interrupt (cause 3) to \code{BASE | 0xc}, four
bytes before \code{BASE + 0xc}. The witness was a nine-frame trace from the
interrupt-entry lemma. Filed as a conformance question, since the documented
alignment rule keeps ordinary software clear of it.

\paragraph{6. The Sail model lets TW be set without U-mode.}
\code{legalize\_mstatus} in the lowRISC revision masks MPRV when U-mode is
absent but not TW, so a \code{csrs} with bit 21 set leaves TW = 1 where Hazard3
correctly hardwires it. The same witness that showed finding 2 showed this
(specification mstatus \code{0x00201888}, core \code{0x00021888}). Upstream
\code{riscv/sail-riscv} fixed it in \code{8de2591} (19 February 2025); the
lowRISC \code{ibex} branch head (\code{5fca3f99}, July 2025) still lacks it.

\paragraph{What the findings are not.}
None is reachable by compiler-generated code; 2 and 3 need non-default
configurations; 5 is documented. They are conformance divergences from the
executable specification, of the kind that testing does not generate cases for
(no test author writes a case for an instruction that should not exist) and
that an equivalence proof produces as a by-product. Finding 4 shows the other
route: a newer conformance suite that enumerates the encodings found it too.

\section{Negative results and lessons}

\begin{itemize}
\item \textbf{Lemma candidates are false about a third of the time at first
  writing, and the false ones are the informative ones.} Fourteen candidates
  were refuted outright. Each refutation was a real behaviour of the core the
  candidate had not allowed for: a load-use stall makes the forwarded operand
  differ from the register file for a cycle; an excepting instruction in the
  M stage leaves a dead slot whose rd matches the next instruction's rs1; a
  late data-phase bus fault turns a locked, non-excepting M stage into an
  excepting one. The operand-tie lemmas went through three versions before the
  guards were right.
\item \textbf{The $k=1$ counterexample does not predict k-induction failure.}
  Sixteen lemmas with a counterexample at $k=1$ proved at depth 1 in the same
  round. A runner that skips the proof on a counterexample loses them.
\item \textbf{Same-layer blindness.} Two lemmas that need each other and sit in
  one layer both time out. They proved at once when placed in consecutive
  layers.
\item \textbf{Stale proofs after a strip.} A goal proven with lemma $L$ beneath
  it stays ``proven'' in the log after $L$ is stripped. The chain checker
  catches it; the fix is to strip and re-submit the dependents fresh, since
  re-inserting a name the checker has seen removed confuses its layer-order
  reconstruction.
\item \textbf{Instrumentation is a constraint.} Probe properties left in the
  script after an experiment are constraints for everything above them, and a
  probe designed to be false makes every proof above it vacuous in the region
  it excludes. Strip them when the experiment ends, and re-prove.
\item \textbf{k-induction with a 400 s cap is the wrong tool for the
  structural lemmas on the full model.} They prove with IC3 in under 30
  minutes and time out under k-induction indefinitely. The model-change
  protocol runs IC3 on those 48 and k-induction on the rest.
\item \textbf{The multiplier does not yield to bit-level SAT.} Bounded checking
  of the unit alone stalls at depth 34; the proof was not attempted further
  and the lemma is an axiom.
\item \textbf{The yosys-slang front end sizes \code{\{32'd0,a\}*\{32'd0,b\}} as
  a 63-bit product.} Explicit 64-bit wires are needed for a checker that
  compares against a full product.
\end{itemize}

\section{Scope and threats to validity}
\label{sec:scope}

\begin{itemize}
\item \textbf{Configuration.} RV32IMC + Zicsr, machine mode only, sequential
  multiplier, no PMP, no debug, no branch predictor. The RP2350 configuration
  is larger in every one of those dimensions. Nothing here is a statement
  about the atomics, the bit-manipulation instructions, U-mode, PMP or the
  debug module.
\item \textbf{The axiom.} The arithmetic of MUL, MULH, MULHSU, MULHU, DIV, DIVU,
  REM and REMU is assumed. A wrong product would not be caught.
\item \textbf{Patched RTL.} The proof is of v1.1.1 with five one-line changes.
  The shipped RTL is not equivalent to the specification; the findings say
  where.
\item \textbf{Coverage narrowings.} Bus faults on either port are reachable
  but not compared: the Sail model has no channel for a bus error, so the
  faulting instruction and its trap are outside the comparison. CSR accesses
  outside the 16-address allowlist drop the ghost until the next query.
  \code{mtval} and \code{mcounteren} are fed as zero and not compared.
\item \textbf{Assumptions.} Hazard3's AHB5 slave assumptions with unbounded
  stalls; debug port tied off (\code{DEBUG\_SUPPORT = 0}); \code{mhartid}
  arbitrary but constant; and the bus returns the data the specification was
  given (\code{MemRdataIsSpecData}, Ibex's \code{mem.sv} assumption), a
  renaming of an otherwise unconstrained free input.
\item \textbf{Not proven.} Liveness (that a query keeps happening), reset
  correctness, and instruction-fetch correctness (that the instruction
  executed at a PC was the one loaded from that PC). All three are open in the
  Ibex proof as well.
\item \textbf{The specification.} A proof against a model is a proof against
  that model. We changed the model in three places; finding 6 shows the model
  can be wrong, and our own patches could be too. Gate 6 restores upstream
  behaviour and was checked to reproduce upstream bit-for-bit at the default
  parameter; Gate 8 matches upstream's later fix; Gate 7 is a configuration
  gate.
\item \textbf{Tool trust.} Yosys, yosys-slang, the Sail back end, psgen,
  aig-manip and rIC3 are all in the trusted base. No proof certificates were
  produced; rIC3 can emit them, and doing so is the obvious next step for
  independent checking.
\item \textbf{Silicon.} The RTL at a commit was verified. Nothing is claimed
  about the RP2350 die.
\end{itemize}

\section{Process}
\label{sec:process}

The harness bring-up, the lemma discovery, the proof runs and the findings were
carried out by an LLM agent (Claude, through Claude Code) driving the tools on
the proof machine, with the author directing. This section records what that
looked like, because it is the part of the method most likely to be reused.

\paragraph{Division of labour.}
The author chose the target, the configuration and the goals; decided each
time whether to fix a divergence in the RTL or to weaken the comparison;
decided to treat the multiplier as an axiom in the Ibex manner rather than
spend days on an algebraic proof; and set the rules the agent ran under. The
agent built the harness gate by gate, wrote every lemma, ran every round, read
every counterexample, wrote the findings and this report's first draft. The
author read the findings against the RTL and the specification text before
accepting them.

\paragraph{Rules that held.}
No goal is reported proven without the chain checker's output, run by hand,
copied verbatim. The author is notified only for a goal proven, a genuine
divergence, or a decision to make. A running proof is never killed to make
room; a round runs to its end. Every result is reported with the model
manifest it came from. A running notebook (\code{RESUME.md}, about 1750
lines by the end) is kept current so that any session can pick up from the
log alone; the notebook is not published, the log is.

\paragraph{What the agent was good at.}
Reading a counterexample to induction and naming the fact that rules it out.
Almost every lemma came from one report: which harness predicates were false
in the final frame, which pipeline signals changed. Turning a refutation into
the right guard (a stall, a dead slot, a late fault) took one to three
attempts. Running the loop unattended for hours, including at night, with the
rules above.

\paragraph{What it was bad at, and what was done about it.}
Declaring victory. Both retractions of 18 September were the agent reporting a
goal as proven on a chain it had not verified; the chain checker exists
because of that, and the rule that no report is made without it. Losing track
of the state after a strip, which is why the runner rather than the agent
does the bookkeeping of what must be re-proven. Killing a stuck proof out of
impatience; the tool permissions were set so that it could not.

\paragraph{Cost.}
The lemma loop took four days of wall-clock time on one machine; the
bring-up before it passed its last three gates on 16 September. The Ibex proof it is modelled on was a multi-month effort by
a specialist; that difference is mostly that the proof structure, the tools
and the shape of every lemma were already known, and that a machine reading
counterexample reports around the clock does not tire.

\section{Related work}

The Ibex proof \cite{ibexformal} is the direct ancestor: the goals, the
follower and ghost, the abstraction, the layered proof script, \code{psgen} and
\code{aig-manip} are theirs, as is the decision to assume the M-type datapath.
It proves about 1000 properties on a 32-core, 128 GB machine, with Jasper or
with the open flow. Hazard3 ships its own riscv-formal \cite{riscvformal}
checks, which are bounded and per-instruction; the FENCE finding was caught
upstream by the riscv-arch-test conformance suite, which enumerates encodings.
The Sail RISC-V model \cite{sail} is the reference in both.

\section{Conclusion}

An RV32IMC, machine-mode configuration of Hazard3 v1.1.1, with five one-line
fixes, is trace equivalent to the RISC-V Sail specification in architectural
state and memory traffic, unboundedly, modulo the arithmetic of its
multiplier/divider and the assumptions of section~\ref{sec:scope}. The proof
used only open tools and one machine for four days of lemma discovery. It
found five conformance divergences in the core and one in the specification
model, all reproduced independently. The multiplier axiom, the larger RP2350
configuration, liveness and proof certificates are the open items, in that
order of difficulty.

\section*{Reproducibility}

The repository holds the harness, the pinned commits and patches, the final
proof script, the three proof logs, the chain verdicts, the tools, and the
findings with their reproducers. \code{reproduce.sh} is meant to rebuild both
models from the pinned sources inside the Ibex formal Nix shell, re-prove
every lemma with nothing stripped, and print the 18 chain verdicts, in about
three hours on 48 cores; the runner scripts still carry the proof machine's
paths and the run has not yet been repeated on a fresh machine
(\code{docs/REPRODUCE.md}). \code{findings/upstream/repro/run.sh} reproduces
the five Hazard3 findings on the CXXRTL testbench in minutes on a laptop.

\begin{thebibliography}{9}
\bibitem{ibexformal} lowRISC and the University of Oxford. \emph{End-to-End
  Formal Verification of Ibex Trace Equivalence against the Sail Model},
  \code{ibex/dv/formal}, Louis-Emile Ploix, 2025.
\bibitem{hazard3} Luke Wren. \emph{Hazard3: a 3-stage RV32IMACZb* processor},
  \url{https://github.com/Wren6991/Hazard3}, v1.1.1.
\bibitem{sail} Alasdair Armstrong et al. \emph{ISA Semantics for ARMv8-A,
  RISC-V, and CHERI-MIPS}, POPL 2019; \code{riscv/sail-riscv}.
\bibitem{riscvformal} Claire Wolf. \emph{riscv-formal},
  \url{https://github.com/YosysHQ/riscv-formal}.
\end{thebibliography}

\end{document}
